\begin{acknowledgements}

论文写至此处，硕士阶段的学习也将告一段落。回顾从本科到研究生的求学经历，老师、同学和家人给予了我许多帮助。在此，向他们表达诚挚的感谢。

首先感谢我的导师王忠杰老师。从本科阶段起，王老师便与我交流生活和成长中的问题，在我面对选择、感到困惑时，耐心听取我的想法，给予坦诚的建议。这些交流让我认识到，学业与生活都需要认真思考，也需要有尝试和承担的勇气。在研究中，老师给予我充分的自由，鼓励我探索感兴趣的方向，并帮助我把握目标。这份信任让我从习惯等待答案、遇到困难容易犹豫，逐渐转向主动分析问题、独立作出判断。老师的教育不仅帮助我完成学业，也改变了我面对生活的态度。今后我会记住这份教诲，保持求知的热情，认真对待自己的选择，以踏实的行动回应老师的信任与期待。

感谢徐晓飞校长。在学习服务计算课程的过程中，我深切感受到了徐老师渊博的知识和严谨的治学态度。课堂上的讲解让我理解服务计算的基本原理，也引导我思考理论与实际需求的联系。老师对问题的认真分析、对学术的热忱以及言传身教，让我受益匪浅，为我此后的学习与研究树立了榜样。

特别感谢贺祥老师。贺老师是我在日常研究中交流最多的老师之一。从本科阶段开展微服务系统治理研究起，老师便帮助我规划时间和研究内容，将宽泛的想法分解为能够逐步完成的工作。研究生期间，每当我在选题、建模和实验中遇到困难，老师都会耐心与我讨论，帮助我找到问题所在，明确下一步的方向。论文写作中，老师投入了大量时间，从内容组织、论证逻辑到文字表述，逐一提出修改意见，帮助我把想法表达清楚。这些细致的指导让我学会了如何推进研究，也让我养成了及时记录、反复检查和认真修改的习惯。感谢老师多年来的耐心与支持，使我在研究受阻时能够理清思路，继续向前。

感谢朝夕相处的师兄师姐和同学们，尤其是王腾、宿梓航、杨智清和云长雨。我们一起探讨学术问题，交流实验进展，也分享生活中的喜悦与忧愁。在彼此的陪伴中，我们共同成长，从同学成为好友。这些相处的点滴，组成了求学期间难忘的时光，也将成为我珍惜的回忆。

感谢父母和家人的理解与支持。你们的信任和包容使我能够安心完成学业，在遇到困难时始终有所依靠。

感谢各位专家在百忙之中抽出时间评审本论文，并提出宝贵意见，谨向各位专家致以最诚挚的谢意。

最后，再次向所有给予我帮助和关注的人表示深深的感谢。
\end{acknowledgements}
